% Источники: exam/materials/преподаватель/2020/23-03-2020-MODEL L5(1).doc; exam/materials/моделирование.pdf, раздел 10.
\section{Стохастическое моделирование. Точечные оценки параметров модели.}

\textbf{Стохастическая модель} — модель, входные данные или параметры которой являются случайными величинами, а преобразование входных данных в выходные осуществляется посредством случайной функции (процесса).

Для определения параметров стохастической модели по выборке $x_1, x_2, \ldots, x_n$ строятся \textbf{точечные оценки}.

\subsection*{Требования к оценкам}

\begin{itemize}
  \item \textbf{Несмещённость:} $M[\hat{\theta}] = \theta$.
  \item \textbf{Состоятельность:} $\hat{\theta} \xrightarrow{P} \theta$ при $n \to \infty$.
  \item \textbf{Эффективность:} минимальная дисперсия среди всех несмещённых оценок.
\end{itemize}

\subsection*{Метод моментов}

Теоретические моменты $\mu_k(\theta)$ приравниваются к выборочным $m_k$:
\[
  \mu_k(\theta) = m_k = \frac{1}{n}\sum_{i=1}^n x_i^k, \quad k = 1, 2, \ldots
\]
Из системы уравнений находят оценки $\hat{\theta}$.

Метод прост, но в общем случае не обеспечивает эффективности.

\subsection*{Метод максимального правдоподобия (ММП)}

Функция правдоподобия:
\[
  L(x_1,\ldots,x_n;\,\theta) = \prod_{i=1}^n f(x_i;\,\theta).
\]
Оценка ММП: $\hat{\theta}$ максимизирует $L$. На практике решают уравнение правдоподобия:
\[
  \frac{\partial \ln L}{\partial \theta} = 0.
\]
ММП-оценка при достаточно общих условиях является \textbf{асимптотически несмещенной} и \textbf{асимптотически эффективной} при $n \to \infty$.

\subsection*{Стандартные выборочные оценки}

\[
  \bar{x} = \frac{1}{n}\sum_{i=1}^n x_i, \qquad
  s^2 = \frac{1}{n-1}\sum_{i=1}^n (x_i - \bar{x})^2.
\]

\clearpage
